\documentclass[10pt,a4paper]{article}
\usepackage{amsmath,amssymb,amsthm,graphicx}
%% Card-local preamble: the source-import-staged retained-preamble.tex of this unit is reproduced verbatim below (author macro and theorem block plus the publisher-class abbreviations it needs).
\usepackage{amsmath,amssymb,amsthm,mathtools}
\usepackage{xcolor}
\usepackage{url}
\usepackage[hidelinks]{hyperref}
\usepackage{xfrac}

\numberwithin{equation}{section}

\newtheorem{Theorem}{Theorem}[section]
\newtheorem*{Theorem*}{Theorem}
\newtheorem{Corollary}[Theorem]{Corollary}
\newtheorem{Lemma}[Theorem]{Lemma}
\newtheorem{prop}[Theorem]{Proposition}
 { \theoremstyle{definition}
\newtheorem{Definition}[Theorem]{Definition}
\newtheorem{Note}[Theorem]{Note}
\newtheorem{Example}[Theorem]{Example}
\newtheorem{Remark}[Theorem]{Remark} }
\newtheorem{rhp}[Theorem]{Riemann--Hilbert Problem}

\usepackage[mathscr]{euscript}
\usepackage{commath, enumitem, tikz}
\newcommand*\up{\ensuremath{\vcenter{\hbox{\includegraphics[width=.6em]{Figures/squareur.pdf}}}}}

\newcommand*\lw{\ensuremath{\vcenter{\hbox{\includegraphics[width=.6em]{Figures/squarell.pdf}}}}}

\newcommand{\ee}{\mathrm{e}}
\newcommand{\dd}{\mathrm{d}}
\newcommand{\ii}{\mathrm{i}}
\newcommand{\para}[1]{\breve{#1}}

\newcommand{\R}{{\mathbb{R}}}
\newcommand{\N}{{\mathbb{N}}}
 \newcommand{\C}{{\mathbb{C}}}
\newcommand{\Q}{{\mathbb{Q}}}
\newcommand{\Z}{{\mathbb{Z}}}
\newcommand\lb{\left(}
\newcommand\rb{\right)}

\newcommand{\re}{\operatorname{Re}}
\newcommand{\im}{\operatorname{Im}}
\newcommand{\Tr}{\operatorname{Tr}}
\newcommand{\rhy}{{\rm RHP}-${\mathbf Y}$}

\newcommand{\ic}{{\mathrm i}}

\newcommand{\supp}{\operatorname{supp}}
\renewcommand{\deg}{\operatorname{deg}}
\renewcommand{\arg}{\operatorname{Arg}}

\newcommand{\sgn}{\operatorname{sgn}}
\newcommand{\qandq}{\qquad \text{and} \qquad}
%\newcommand{\qasq}{\quad \text{as} \quad }
%\newcommand{\qforq}{\quad \text{for} \quad }
%\newcommand{\br}{\textcolor{brown} }

\usepackage{cleveref}
\crefname{eq}{original equation}{original equations}
\Crefname{eq}{Original equation}{Original equations}
\makeatletter
\newlabel{eq:$D_6$-parameters}{{2}{source}{}{sourcepointer-1}{}}
\newlabel{eq:$D_6$-parameters@cref}{{[eq][0][]2}{[1][0][]source}{}{}{}}
\makeatother
\crefname{eq}{original equation}{original equations}
\Crefname{eq}{Original equation}{Original equations}
\makeatletter
\newlabel{eq:PIII-$D_8$}{{1}{source}{}{sourcepointer-2}{}}
\newlabel{eq:PIII-$D_8$@cref}{{[eq][0][]1}{[1][0][]source}{}{}{}}
\makeatother
\crefname{eq}{original equation}{original equations}
\Crefname{eq}{Original equation}{Original equations}
\makeatletter
\newlabel{eq:V-coeffs}{{2}{source}{}{sourcepointer-3}{}}
\newlabel{eq:V-coeffs@cref}{{[eq][0][]2}{[1][0][]source}{}{}{}}
\makeatother
\crefname{eq}{original equation}{original equations}
\Crefname{eq}{Original equation}{Original equations}
\makeatletter
\newlabel{eq:V-coeffsU1}{{2}{source}{}{sourcepointer-4}{}}
\newlabel{eq:V-coeffsU1@cref}{{[eq][0][]2}{[1][0][]source}{}{}{}}
\makeatother
\crefname{eq}{original equation}{original equations}
\Crefname{eq}{Original equation}{Original equations}
\makeatletter
\newlabel{eq:c1-exact}{{2}{source}{}{sourcepointer-5}{}}
\newlabel{eq:c1-exact@cref}{{[eq][0][]2}{[1][0][]source}{}{}{}}
\makeatother
\crefname{eq}{original equation}{original equations}
\Crefname{eq}{Original equation}{Original equations}
\makeatletter
\newlabel{eq:v-ODE}{{2}{source}{}{sourcepointer-6}{}}
\newlabel{eq:v-ODE@cref}{{[eq][0][]2}{[1][0][]source}{}{}{}}
\makeatother
\crefname{eq}{original equation}{original equations}
\Crefname{eq}{Original equation}{Original equations}
\makeatletter
\newlabel{eq:v-series}{{2}{source}{}{sourcepointer-7}{}}
\newlabel{eq:v-series@cref}{{[eq][0][]2}{[1][0][]source}{}{}{}}
\makeatother
\crefname{eq}{original equation}{original equations}
\Crefname{eq}{Original equation}{Original equations}
\makeatletter
\newlabel{eq:xi-general-exact}{{2}{source}{}{sourcepointer-8}{}}
\newlabel{eq:xi-general-exact@cref}{{[eq][0][]2}{[1][0][]source}{}{}{}}
\makeatother
\crefname{eq}{original equation}{original equations}
\Crefname{eq}{Original equation}{Original equations}
\makeatletter
\newlabel{eq:xi-inequalities}{{2}{source}{}{sourcepointer-9}{}}
\newlabel{eq:xi-inequalities@cref}{{[eq][0][]2}{[1][0][]source}{}{}{}}
\makeatother
\title{The D6-D8 series limit and its Painleve III(D8) limit}
\author{Barhoumi, Ahmad, Lisovyy, Oleg, Miller, Peter D., Prokhorov, Andrei}
\date{doi:10.3842/SIGMA.2024.019; CC BY 4.0}
\begin{document}
\maketitle
\setcounter{section}{2}
\setcounter{Theorem}{0}
\noindent\textit{Editorial note.} All delivered mathematical text is the author's own, copied from the pinned journal source file sources/sigma-2024-019-painlev-233-iii-monodromy-maps-under-the-d-6/pinned-original-source.tex (Painlevé-III Monodromy Maps Under the D6 to D8 Confluence and Applications to the Large-Parameter Asymptotics of Rational Solutions); no statement or proof step was rewritten, reordered or removed.\par\medskip
\noindent\textit{Editorial note.} The counted proof is the author's complete argument: All delivered mathematical text is the author own, copied verbatim from the pinned journal source. Statement source lines 471-477 with the author own proof at 478-481 (label thm:\$D\_6\$-\$D\_8\$-series). Printed number 2.1 (the census published value, independently reproduced by the source walk: counter 0 of section 2). The statement and proof spans touch neither the public statement [255,284] nor the public proof boxes [1767,2887], [3074,3310] of problem\_id 902, and the spans carry no unstaged citation.\par\medskip
\noindent\textit{Editorial note.} The label thm:\$D\_6\$-\$D\_8\$-series names the counted statement, and the proof environment immediately following it in the source is delivered verbatim as its proof.\par\medskip
\noindent\textit{Editorial note.} The counted statement is carried by the author's own Theorem environment in the article source.\par\medskip
\noindent\textit{Editorial note.} Journal-unit assembly: card-local preamble (author theorem environments and macros verbatim), article class, bibliography backend none; any literature citation carried by the delivered span is disclosed in the notes of this card.\par\medskip
\noindent\textit{Editorial note.} Revision r002: re-assembled once the delivery chain admitted "\$" to the external-label whitelist (Chief ruling §8.8.136, 2026-10-15 01:5x), which two pointer labels of this card need (eq:\$D\_6\$-parameters, eq:PIII-\$D\_8\$). Nothing else changed.\par\medskip
\noindent\textit{Editorial scope.} the counted statement and its own proof environment only\par\medskip
\noindent\textit{Allowed standard prerequisites (editorial).}\begin{itemize}
\item The sequence of solutions U\_n. Supplies background that the delivered proof uses. The entry is declared as a dependency; it is not re-proved on the card, and the load-bearing steps of the counted statement are carried out inside the delivered span.
\end{itemize}\medskip
\ifdefined\hypertarget\hypertarget{sourcepointer-1}{}\fi
\noindent\textit{Source reference.} Original equation 2: operations/codex-takeover-01a12558/corpus-repair-20261010/research-line/harness-ready-20261011/sources/sigma-2024-019-painlev-233-iii-monodromy-maps-under-the-d-6/pinned-original-source.tex, source lines 370--370. Navigation pointer invoked at line(s) 480: the counted text refers to this environment of the same journal unit, which is not retained in this card, so the reference is declared as a pointer carrying the number the article itself prints for it (the section and the place of the statement inside that section, derived from the shared section-scoped theorem counter of the source); the cited item is an allowed input and none of it is reproduced inside the span.\par\medskip
\ifdefined\hypertarget\hypertarget{sourcepointer-2}{}\fi
\noindent\textit{Source reference.} Original equation 1: operations/codex-takeover-01a12558/corpus-repair-20261010/research-line/harness-ready-20261011/sources/sigma-2024-019-painlev-233-iii-monodromy-maps-under-the-d-6/pinned-original-source.tex, source lines 131--131. Navigation pointer invoked at line(s) 475: the counted text refers to this environment of the same journal unit, which is not retained in this card, so the reference is declared as a pointer carrying the number the article itself prints for it (the section and the place of the statement inside that section, derived from the shared section-scoped theorem counter of the source); the cited item is an allowed input and none of it is reproduced inside the span.\par\medskip
\ifdefined\hypertarget\hypertarget{sourcepointer-3}{}\fi
\noindent\textit{Source reference.} Original equation 2: operations/codex-takeover-01a12558/corpus-repair-20261010/research-line/harness-ready-20261011/sources/sigma-2024-019-painlev-233-iii-monodromy-maps-under-the-d-6/pinned-original-source.tex, source lines 426--426. Navigation pointer invoked at line(s) 479: the counted text refers to this environment of the same journal unit, which is not retained in this card, so the reference is declared as a pointer carrying the number the article itself prints for it (the section and the place of the statement inside that section, derived from the shared section-scoped theorem counter of the source); the cited item is an allowed input and none of it is reproduced inside the span.\par\medskip
\ifdefined\hypertarget\hypertarget{sourcepointer-4}{}\fi
\noindent\textit{Source reference.} Original equation 2: operations/codex-takeover-01a12558/corpus-repair-20261010/research-line/harness-ready-20261011/sources/sigma-2024-019-painlev-233-iii-monodromy-maps-under-the-d-6/pinned-original-source.tex, source lines 422--422. Navigation pointer invoked at line(s) 479: the counted text refers to this environment of the same journal unit, which is not retained in this card, so the reference is declared as a pointer carrying the number the article itself prints for it (the section and the place of the statement inside that section, derived from the shared section-scoped theorem counter of the source); the cited item is an allowed input and none of it is reproduced inside the span.\par\medskip
\ifdefined\hypertarget\hypertarget{sourcepointer-5}{}\fi
\noindent\textit{Source reference.} Original equation 2: operations/codex-takeover-01a12558/corpus-repair-20261010/research-line/harness-ready-20261011/sources/sigma-2024-019-painlev-233-iii-monodromy-maps-under-the-d-6/pinned-original-source.tex, source lines 390--390. Navigation pointer invoked at line(s) 480: the counted text refers to this environment of the same journal unit, which is not retained in this card, so the reference is declared as a pointer carrying the number the article itself prints for it (the section and the place of the statement inside that section, derived from the shared section-scoped theorem counter of the source); the cited item is an allowed input and none of it is reproduced inside the span.\par\medskip
\ifdefined\hypertarget\hypertarget{sourcepointer-6}{}\fi
\noindent\textit{Source reference.} Original equation 2: operations/codex-takeover-01a12558/corpus-repair-20261010/research-line/harness-ready-20261011/sources/sigma-2024-019-painlev-233-iii-monodromy-maps-under-the-d-6/pinned-original-source.tex, source lines 365--365. Navigation pointer invoked at line(s) 472: the counted text refers to this environment of the same journal unit, which is not retained in this card, so the reference is declared as a pointer carrying the number the article itself prints for it (the section and the place of the statement inside that section, derived from the shared section-scoped theorem counter of the source); the cited item is an allowed input and none of it is reproduced inside the span.\par\medskip
\ifdefined\hypertarget\hypertarget{sourcepointer-7}{}\fi
\noindent\textit{Source reference.} Original equation 2: operations/codex-takeover-01a12558/corpus-repair-20261010/research-line/harness-ready-20261011/sources/sigma-2024-019-painlev-233-iii-monodromy-maps-under-the-d-6/pinned-original-source.tex, source lines 385--385. Navigation pointer invoked at line(s) 479: the counted text refers to this environment of the same journal unit, which is not retained in this card, so the reference is declared as a pointer carrying the number the article itself prints for it (the section and the place of the statement inside that section, derived from the shared section-scoped theorem counter of the source); the cited item is an allowed input and none of it is reproduced inside the span.\par\medskip
\ifdefined\hypertarget\hypertarget{sourcepointer-8}{}\fi
\noindent\textit{Source reference.} Original equation 2: operations/codex-takeover-01a12558/corpus-repair-20261010/research-line/harness-ready-20261011/sources/sigma-2024-019-painlev-233-iii-monodromy-maps-under-the-d-6/pinned-original-source.tex, source lines 403--403. Navigation pointer invoked at line(s) 480: the counted text refers to this environment of the same journal unit, which is not retained in this card, so the reference is declared as a pointer carrying the number the article itself prints for it (the section and the place of the statement inside that section, derived from the shared section-scoped theorem counter of the source); the cited item is an allowed input and none of it is reproduced inside the span.\par\medskip
\ifdefined\hypertarget\hypertarget{sourcepointer-9}{}\fi
\noindent\textit{Source reference.} Original equation 2: operations/codex-takeover-01a12558/corpus-repair-20261010/research-line/harness-ready-20261011/sources/sigma-2024-019-painlev-233-iii-monodromy-maps-under-the-d-6/pinned-original-source.tex, source lines 418--418. Navigation pointer invoked at line(s) 479: the counted text refers to this environment of the same journal unit, which is not retained in this card, so the reference is declared as a pointer carrying the number the article itself prints for it (the section and the place of the statement inside that section, derived from the shared section-scoped theorem counter of the source); the cited item is an allowed input and none of it is reproduced inside the span.\par\medskip
\begin{Theorem}\label{thm:$D_6$-$D_8$-series}
Fix $\alpha\in\mathbb{C}$ and $\beta\in\mathbb{C}$ and let $\{U_n(z;\alpha,\beta)\}_{n=1}^\infty$ be a sequence of solutions of~\eqref{eq:v-ODE} that are analytic at the origin $z=0$ and suppose that
\[\lim_{n\to\infty}U_n(0;\alpha,\beta)=\upsilon_{\infty,0}=\upsilon_{\infty,0}(\alpha,\beta)\neq 0.\]
 Then there exists a radius $\rho>0$ such that for all $n$ sufficiently large $U_n(z;\alpha,\beta)$ is analytic for $|z|<\rho$ and such that $U_n(z;\alpha,\beta)\to U_\infty(z;\alpha,\beta)$ as $n\to\infty$ uniformly for $|z|<\rho$,
where~${U(z)=U_\infty(z;\alpha,\beta)}$ is the unique solution of the Painlev\'e-{\rm III}$(D_8)$ equation \eqref{eq:PIII-$D_8$}
that is analytic at the origin with $U_\infty(0;\alpha,\beta)=\upsilon_{\infty,0}$.
\end{Theorem}

\begin{proof}
Let $\{\upsilon_{n,k}\}_{k=0}^\infty$ denote the power series coefficients of $U_n(z;\alpha,\beta)$ as in \eqref{eq:v-series}. Define~$\Upsilon_0$ by~$\Upsilon_0>2|\upsilon_{\infty,0}|$ (say), and obtain the subsequent coefficients $\{\Upsilon_k\}_{k=1}^\infty$ via \eqref{eq:V-coeffsU1}--\eqref{eq:V-coeffs}. Comparing \eqref{eq:xi-inequalities}--\eqref{eq:V-coeffs} then shows that for all $n$ sufficiently large,
$|\upsilon_{n,k}|\le \Upsilon_k$ holds for all \linebreak ${k=0,1,2,\dots}$. For each fixed $k=0,1,2,\dots$, the recurrence relations \eqref{eq:c1-exact}--\eqref{eq:xi-general-exact} together with the limit ${\upsilon_{n,0}\to\upsilon_{\infty,0}}$ show that $\upsilon_{n,k}$ tends to a limiting value $\upsilon_{\infty,k}$ as $n\to\infty$, with $|\upsilon_{\infty,k}|\le\Upsilon_k$. The convergence of $U_n(z;\alpha,\beta)$ to a limiting analytic function $U_\infty(z;\alpha,\beta)$ with $U_\infty(0;\alpha,\beta)=\upsilon_{\infty,0}(\alpha,\beta)$ then follows by dominated convergence. That the limiting analytic function $U_\infty(z;\alpha,\beta)$ is a~solution of~\eqref{eq:PIII-$D_8$} follows from passing to the limit in each term of \eqref{eq:v-ODE} using \eqref{eq:$D_6$-parameters}. That this solution is the unique analytic solution of \eqref{eq:PIII-$D_8$} with the specified value at $z=0$ then follows from passing to the limit in the recurrence relations \eqref{eq:c1-exact}--\eqref{eq:xi-general-exact}.
\end{proof}

\end{document}
